% Figure 6 -- the accuracy / harm trade-off. Both panels use zero-based axes.
% Data: figures/data/fig6_accuracy_harm.dat, fig6b_conflict_handling.dat
% Generated by PHASE5_PUBLICATION/figures_source/make_figure_data.py from
% PHASE3/results/per_model_config.csv (primary models) and PHASE3/ANALYSIS/METRICS.json
\begin{figure}[t]
\centering
\begin{tikzpicture}
% ---------------- panel (a): success and harm side by side, both measured from zero
\begin{axis}[
  paperaxis, name=panelA,
  width=0.40\linewidth, height=7.0cm,
  xbar=0pt, bar width=2.4pt,
  xmin=0, xmax=78,
  xlabel={\% of runs},
  symbolic y coords={EBG-noR3,EBG-unsafe,EBG-all,EBG,EB,LG,LA,RT,R,T,G,B},
  ytick=data,
  y tick label style={font=\scriptsize},
  xmajorgrids, xtick={0,20,40,60},
  enlarge y limits=0.05,
  legend style={at={(0.5,-0.19)}, anchor=north, legend columns=2, draw=cNeutral},
  legend cell align=left,
  title={(a) success and harm, from zero},
]
\addplot[draw=cAcc, fill=cAcc, fill opacity=0.85]
  table[col sep=tab, x=success, y=config] {figures/data/fig6_accuracy_harm.dat};
\addplot[draw=cHarm, fill=cHarm, fill opacity=0.85, postaction={pattern=north east lines, pattern color=white}]
  table[col sep=tab, x=harmful, y=config] {figures/data/fig6_accuracy_harm.dat};
\legend{task success, harmful answer}
\end{axis}

% ---------------- panel (b): the trade-off plane
\begin{axis}[
  paperaxis, name=panelB, at={(panelA.right of east)}, anchor=left of west, xshift=8mm,
  width=0.47\linewidth, height=7.0cm,
  xlabel={task success (\% of runs) $\rightarrow$ better},
  ylabel={harmful answers (\% of runs) $\rightarrow$ worse},
  xmin=0, xmax=82, ymin=0, ymax=26,
  xtick={0,20,40,60,80}, ytick={0,5,10,15,20,25},
  xmajorgrids, ymajorgrids,
  title={(b) the twelve conditions in the trade-off plane},
  clip=false,
]
% the trade we draw attention to: gating gives up accuracy to buy safety
\draw[->, >=latex, line width=1.1pt, draw=cAll] (axis cs:61.3,3.1) -- (axis cs:71.1,6.9);
\draw[draw=cAll, line width=0.4pt] (axis cs:67,5.1) -- (axis cs:62,13.5);
\node[lbl, text=cAll, anchor=south, align=center] at (axis cs:62,13.8)
  {\scriptsize \textbf{drop the gate}\\[-1pt]\scriptsize $+9.7$\pp{} success\\[-1pt]\scriptsize $+3.8$\pp{} harm};

\addplot[only marks, mark=*, mark size=2.0pt, draw=cNone, fill=cNone] coordinates {(46.0,7.0)};
\node[lbl, anchor=north east] at (axis cs:46.0,7.0) {B\,};
\addplot[only marks, mark=square*, mark size=1.9pt, draw=cAdv!65!black, fill=cAdv]
  coordinates {(43.5,7.9) (46.9,7.4)};
\node[lbl, anchor=south east] at (axis cs:43.5,7.9) {G\,};
\node[lbl, anchor=west] at (axis cs:46.9,7.4) {\,T};
\addplot[only marks, mark=triangle*, mark size=2.4pt, draw=cEnf, fill=cEnf]
  coordinates {(45.0,16.1) (44.2,22.6)};
\node[lbl, anchor=east] at (axis cs:45.0,16.1) {R\,};
\node[lbl, anchor=east] at (axis cs:44.2,22.6) {RT\,};
\addplot[only marks, mark=diamond*, mark size=2.4pt, draw=cLLM!70!black, fill=cLLM]
  coordinates {(47.9,13.2) (45.4,21.3)};
\node[lbl, anchor=west] at (axis cs:47.9,13.2) {\,LA};
\node[lbl, anchor=west] at (axis cs:45.4,21.3) {\,LG};
\addplot[only marks, mark=pentagon*, mark size=2.2pt, draw=cEvarb, fill=cEvarb]
  coordinates {(55.3,2.7) (61.5,3.3) (60.0,3.4)};
\node[lbl, anchor=south] at (axis cs:55.3,2.9) {EB};
\addplot[only marks, mark=*, mark size=3.0pt, draw=cEBG, fill=cEBG] coordinates {(61.3,3.1)};
\node[lbl, anchor=north, text=cEBG] at (axis cs:60.6,2.6) {\textbf{EBG}};
\addplot[only marks, mark=*, mark size=3.0pt, draw=cAll, fill=cAll] coordinates {(71.1,6.9)};
\node[lbl, anchor=south east, text=cAll] at (axis cs:71.8,7.8) {\textbf{EBG-all}};
\end{axis}
\end{tikzpicture}

\vspace{1mm}
\caption[Accuracy and harm are not improved together]{%
\textbf{Accuracy and harm do not move together.} Pooled over the four primary models and four decodings,
$n=4288$ runs per condition. \textbf{(a)} Task success and harmful-answer rate for all twelve conditions, both
measured from zero on a common scale. Enforcing record-derived labels (\cond{RT}) roughly triples the harm rate
relative to plain memory (22.6\% against 7.0\%) without improving success. \textbf{(b)} The same conditions placed
by success and harm. \cond{EBG-unsafe} and \cond{EBG-noR3} sit within half a point of \cond{EBG} and are drawn in
the same green but left unlabelled to keep the cluster legible; their exact rates are in
\Cref{tab:rates}. The amber arrow is the central counter-result of this paper: removing the evidence gate, so
that \emph{every} arbitration is enforced rather than only strong-strength ones, moves the system to
\cond{EBG-all} --- $9.7$ percentage points more task success and $3.8$ more harm, with arbitrary picks rising from
13 runs to 166 and false resolution from 1.9\% to 12.9\%. Gating is a trade, not a dominance.\\
\emph{Accessible description:} in panel (b) the evidence-basis conditions occupy the lower right (high success,
low harm); the enforced and LLM-arbiter conditions occupy the upper left (no better success, three to seven times
the harm); an arrow runs up and to the right from \cond{EBG} to \cond{EBG-all}.\\
\emph{Source:} \srcfile{PHASE3/results/per_model_config.csv}; plotted from
\srcfile{figures/data/fig6_accuracy_harm.dat}.}
\label{fig:tradeoff}
\end{figure}
